\section*{Scope}
\label{sec:scope}

This page is the paper in brief. Each line names what it rests on, and the sections after give the numbers.

\subsection*{Filled in since \emph{The Missing Amplitude}}

\begin{itemize}
\item \textbf{The missing amplitude is supplied.} The rule now carries exact amplitudes in $\mathbb{Z}[\omega][1/2]$.
The meeting (the fear beat) is the one non-Clifford gate, and a moving vibe is quantum while a static one is classical
(Section~\ref{sec:light}).
\item \textbf{A rule with a reason for every piece}: the $W(F_4)$ knit, the collision, the doublet lock, the meeting,
the pass, the no-veto store and the coin, each chosen over a stated alternative (Section~\ref{sec:rule}).
\item \textbf{A working vacuum}, balanced, reversible and covariant on 51 of 51 cases (Section~\ref{sec:rule}).
\item \textbf{Light}, as compact $U(1)$ from column sums, with two polarizations and Coulomb's $1/(24\pi r)$
(Section~\ref{sec:light}).
\item \textbf{Quantum results on the current rule}: Bell, Tsirelson, contextuality as fear, no-cloning, teleportation,
GHZ, spin-statistics, Pauli exclusion, and the second law with its arrow (Section~\ref{sec:light}).
\end{itemize}

\subsection*{Solved}

\begin{itemize}
\item \textbf{Gravity as a mechanism.} One trit per link, summed into a depth, gives $1/r$, the same fall for
everything larger than a dock, light bending by 2, a pull at $c$, and $\alpha$ held constant (Section~\ref{sec:gravity}).
\item \textbf{Horizons}: a radius of $2GM$, Hawking's redshift scaling, particles from a forming horizon, and an area
law with Bekenstein's bound as dynamics (Section~\ref{sec:gravity}).
\item \textbf{Spin 2 derived, not chosen}: sliding the unlabeled docks forces linearized Einstein--Hilbert, the ADM
structure, a wave speed of exactly $c$ and $\lambda = 1$ (Section~\ref{sec:gravity}).
\item \textbf{Randall--Sundrum's short-range number} in the flat bulk, with its 4/3 factor derived
(Section~\ref{sec:gravity}).
\item \textbf{A composite at relativistic speed} on a line, and falling alike (Section~\ref{sec:changelog}).
\item \textbf{Motion in 3d, one-body.} Line-only motion is ruled out by observation. A stated rule change gives every
excitation one limiting speed $c^* = c/2$ and exact relativistic, isotropic motion, in the many-body rule, with the
vacuum inert and gravity's sign kept (Section~\ref{sec:motion}).
\item \textbf{The binding problem, narrowed by theorems}: the band sum, the Kesten floor, theorems A and B, the first
order walls of $\mathbb{Z}_3$ and of the bulk, and a $2I$ string that loosens continuously on the husk
(Section~\ref{sec:binding}).
\end{itemize}

\subsection*{Still open}

\begin{itemize}
\item \textbf{Binding that moves.} Whether the beat keeps the ordered $2I$ vacuum long enough, and then a light bound
member moving in 3d with inertia equal to energy. This is the program's sharpest open question.
\item \textbf{A Planck temperature} for horizons, and the bulk's warp, radion and brane bending.
\item \textbf{One light speed} for the husk light on the candidate rule.
\item \textbf{The electron across lines and $g$}, \textbf{the weak force and CP violation}, \textbf{the value of
$\alpha$}, \textbf{a bounded self}, and \textbf{one rule for every row} (Section~\ref{sec:open}).
\end{itemize}

The candidate rule of Section~\ref{sec:motion} is not yet the working rule: it is adopted only once binding works
on it.
